\section{Файл заголовков Resource.h}
\footnotesize
\begin{listing}{1}
/*************************************************************************
* КАФЕДРА № 304 | 2 КУРС | УЧЕБНАЯ ПРАКТИКА                             *
*------------------------------------------------------------------------*
* Project Name  : Программа расчета параметров векторных полей           *
* (Теория поля)                                                          *
* File Name     : Resource.h                                             *
* Author(s)     : Варфоломеев К.В., Нагороднюк А.Д.,                     *
*                 Великоцкая А.С.                                        *
* Created       : 02/07/2026                                             *
* Last Modified : 09/07/2026                                             *
* Version       : 1.0                                                    *
* Copyright     : (c) 2026. All rights reserved.                         *
* License       : Educational use only                                   *
*------------------------------------------------------------------------*
* НАЗНАЧЕНИЕ ФАЙЛА:                                                      *
* Содержит макроопределения идентификаторов графических ресурсов и       *
* элементов управления пользовательского интерфейса Win32.               *
*                                                                        *
* ОПИСАНИЕ ФУНКЦИОНИРОВАНИЯ:                                             *
* Файл Resource.h автоматически генерируется средой разработки           *
* Microsoft Visual Studio и дополняется разработчиком для связывания     *
* визуальных компонентов (кнопок, диалоговых окон, иконок) с их          *
* числовыми идентификаторами в таблице ресурсов приложения.              *
*                                                                        *
* В ДАННОМ ФАЙЛЕ ОПРЕДЕЛЕНЫ:                                             *
* 1. Системные идентификаторы приложения (IDS_APP_TITLE,                 *
* IDR_MAINFRAME, IDI_VECTORFIELD, IDI_SMALL).                            *
* 2. Идентификаторы диалоговых окон меню и подсистемы "О программе"      *
* (IDD_VECTORFIELD_DIALOG, IDD_ABOUTBOX).                                *
* 3. Уникальные идентификаторы управляющих кнопок на панели              *
* инструментов, обеспечивающие обработку пользовательских событий        *
* в оконной процедуре:                                                   *
* - ID_BTN_IMPORT : Импорт данных из внешнего файла .txt                 *
* - ID_BTN_SHOW   : Отображение матрицы полиномов                        *
* - ID_BTN_CALC   : Запуск процедур математического анализа полей        *
* - ID_BTN_EXPORT : Экспорт сформированного отчета на диск               *
*************************************************************************/

//{{NO_DEPENDENCIES}}
// Включаемый файл, созданный в Microsoft Visual C++.
// Используется VectorField.rc
#define IDS_APP_TITLE           103
#define IDR_MAINFRAME           128
#define IDD_VECTORFIELD_DIALOG  102
#define IDD_ABOUTBOX            103
#define IDM_ABOUT               104
#define IDM_EXIT                105
#define IDI_VECTORFIELD         107
#define IDI_SMALL               108
#define IDC_VECTORFIELD         109
#define IDC_MYICON              2
#ifndef IDC_STATIC
#define IDC_STATIC              -1
#endif
// Кнопки интерфейса
#define ID_BTN_IMPORT           1001
#define ID_BTN_SHOW             1002
#define ID_BTN_CALC             1003
#define ID_BTN_VIS_A            1004
#define ID_BTN_VIS_B            1005
#define ID_BTN_EXPORT           1006
// Следующие стандартные значения для новых объектов
#ifdef APSTUDIO_INVOKED
#ifndef APSTUDIO_READONLY_SYMBOLS
#define _APS_NO_MFC             130
#define _APS_NEXT_RESOURCE_VALUE 129
#define _APS_NEXT_COMMAND_VALUE 32771
#define _APS_NEXT_CONTROL_VALUE 1007
#define _APS_NEXT_SYMED_VALUE   110
#endif
#endif
\end{listing}
\normalsize

\pagebreak

\section{Файл заголовков VectorField.h}
\footnotesize
\begin{listing}{1}
/*************************************************************************
* КАФЕДРА № 304 | 2 КУРС | УЧЕБНАЯ ПРАКТИКА                             *
*------------------------------------------------------------------------*
* Project Name  : Программа расчета параметров векторных полей           *
* (Теория поля)                                                          *
* File Name     : VectorField.h                                          *
* Author(s)     : Варфоломеев К.В., Нагороднюк А.Д.,                     *
*                 Великоцкая А.С.                                        *
* Created       : 02/07/2026                                             *
* Last Modified : 09/07/2026                                             *
* Version       : 1.0                                                    *
* Copyright     : (c) 2026. All rights reserved.                         *
* License       : Educational use only                                   *
*------------------------------------------------------------------------*
* НАЗНАЧЕНИЕ ФАЙЛА:                                                      *
* Содержит объявления базовых структур данных трехмерного пространства,  *
* класса символьных полиномов и статических методов аналитического       *
* ядра для вычислений в теории поля.                                     *
*                                                                        *
* ОПИСАНИЕ ИНТЕРФЕЙСОВ И ТИПОВ ДАННЫХ:                                   *
* В заголовочном файле сосредоточены все типы данных, на которых         *
* базируется вычислительная логика программы:                            *
*                                                                        *
* 1. Структура Vector3D:                                                 *
* Описывает трехмерный геометрический вектор с вещественными             *
* координатами (x, y, z).                                                *
*                                                                        *
* 2. Структура Monomial:                                                 *
* Описывает структуру отдельного монома, содержащего вещественный        *
* коэффициент "coef" и целочисленные степени трех пространственных       *
* переменных (px, py, pz).                                               *
*                                                                        *
* 3. Класс Polynomial:                                                   *
* Инкапсулирует логику символьной работы со многочленами (сложение,      *
* зануление переменных при расчете потенциалов, символьное               *
* интегрирование, численное вычисление значения в точке).                *
*                                                                        *
* 4. Класс FieldMathCalculations:                                        *
* Математическое ядро программы, выполняющее синтаксический анализ       *
* строковых формул, расчет дивергенции, ротора, потока (через теорему    *
* Гаусса), циркуляции (через теорему Стокса), а также нахождение         *
* потенциала векторного поля.                                            *
*************************************************************************/
#pragma once
#include <string>
#include <vector>
#include <functional>
#include "resource.h"
using namespace std;
// =======================================================================
// ОБЪЯВЛЕНИЯ СТРУКТУР И КЛАССОВ
// =======================================================================
struct Vector3D {double x, y, z;}; struct Monomial {double coef; int px, py, pz;};
class Polynomial {
public:
    vector<Monomial> terms; Polynomial(initializer_list<Monomial> init);
    Polynomial(); Polynomial SetVariableToZero(bool zeroX, bool zeroY,
    bool zeroZ) const;
    Polynomial IntegrateSymbolically(int varIndex) const;
    void Add(const Polynomial& other); wstring ToWString() const;
    double Evaluate(double x, double y, double z) const;
};
class FieldMathCalculations {
public: typedef function<double(double, double, double)> FUNC3D;
    typedef function<Vector3D(double, double, double)> FUNC3D_VEC;
    static Polynomial fieldA_X, fieldA_Y, fieldA_Z;
    static Polynomial fieldB_X, fieldB_Y, fieldB_Z;
    static void ParsePolynomials(const wstring& a_str, const wstring& b_str);
    static Vector3D EvaluateFieldA(double x, double y, double z);
    static Polynomial GetFieldBX(); static Polynomial GetFieldBY();
    static Polynomial GetFieldBZ();
    static Vector3D EvaluateFieldB(double x, double y, double z);
    static double Calc_PartialDerivative(FUNC3D f, int variable, 
                                       double x, double y, double z);
    static double IntegrateQuarterDisk(function<double(double, double)> func, 
                                       double radius = 1.0, int steps = 50);
    static double Calc_Divergence(FUNC3D_VEC field, double x, double y, double z);
    static Vector3D Calc_Rotor(FUNC3D_VEC field, double x, double y, double z);
    static wstring DetermineFieldClass(FUNC3D_VEC field);
    static double Calc_Flux_GaussOstrogradsky();
    static double Calc_Flux_Direct(); static double Calc_Circulation_Stokes();
    static double Calc_Circulation_Direct(); static Polynomial GetAnalyticalPotential();
    static bool VerifyPotentialCorrectness(); static bool CheckRotorBZero();
private: static Polynomial ParsePoly(wstring expr);
    static void ExtractComponents(wstring line, Polynomial& P, 
                                  Polynomial& Q, Polynomial& R);
    static void ReplaceAll(wstring& str, const wstring& from, const wstring& to);
};
\end{listing}
\normalsize

\pagebreak

\section{Метод Polynomial::SetVariableToZero}
\footnotesize
\begin{listing}{1}
/*************************************************************************
* КАФЕДРА № 304 | 2 КУРС | УЧЕБНАЯ ПРАКТИКА                              *
*------------------------------------------------------------------------*
* Проект        : Программа расчета параметров векторных полей           *
* Файл          : VectorField.cpp                                        *
* Создан        : 02.07.2026                                             *
*                                                                        *
* РАСПРЕДЕЛЕНИЕ ОТВЕТСТВЕННОСТИ (ФУНКЦИОНАЛЬНЫЕ МОДУЛИ):                 *
* - Великоцкая А.С. : Модули 1-20 (Математическое ядро)                  *
* - Варфоломеев К.В. : Модули 20-35 (Вычисления и интерфейс)             *
* - Нагороднюк А.Д. : Модули 35-49 (Интерфейс и визуализация)            *
*------------------------------------------------------------------------*
* НАЗНАЧЕНИЕ ФАЙЛА:                                                      *
* Реализация логики работы системы: описание полей, алгоритмы анализа,   *
* обработка данных и управление интерфейсом (WinAPI).                    *
*************************************************************************/
/*************************************************************************
* Метод         : Polynomial::SetVariableToZero                          *
*------------------------------------------------------------------------*
* Назначение    : Символьное зануление переменных в полиноме             *
* Входные       : zeroX, zeroY, zeroZ - флаги зануления осей x, y, z     *
* Выходные      : Новый объект Polynomial с отфильтрованными мономами    *
*------------------------------------------------------------------------*
* ОПИСАНИЕ ФУНКЦИОНИРОВАНИЯ:                                             *
* Метод осуществляет символьное зануление выбранных переменных в         *
* структуре полинома. Это необходимо для подготовки к аналитическому     *
* интегрированию (потенциала поля). Метод фильтрует мономы, отбрасывая   *
* те, у которых степени зануляемых переменных больше нуля.               *
*************************************************************************/
Polynomial Polynomial::SetVariableToZero(bool zeroX, bool zeroY, 
                                         bool zeroZ) const {
    Polynomial result;
    for (const auto& term : terms) {
        if ((zeroX && term.px > 0) || (zeroY && term.py > 0) || 
            (zeroZ && term.pz > 0)) continue;
        result.terms.push_back(term);
    }
    return result;
}
\end{listing}
\normalsize

\pagebreak

\section{Метод Polynomial::IntegrateSymbolically}
\footnotesize
\begin{listing}{1}
/*************************************************************************
* Метод         : Polynomial::IntegrateSymbolically                      *
*------------------------------------------------------------------------*
* Назначение    : Символьное интегрирование полинома по оси              *
* Входные       : varIndex - индекс переменной (0 -> x, 1 -> y, 2 -> z)  *
* Выходные      : Новый объект Polynomial (первообразная)                *
*------------------------------------------------------------------------*
* ОПИСАНИЕ ФУНКЦИОНИРОВАНИЯ:                                             *
* Выполняет аналитическое интегрирование полинома по выбранной оси.      *
* Метод обходит каждый моном и применяет правило интегрирования для      *
* степенной функции. Результат — новая первообразная функция,            *
* используемая для точного расчета потенциала.                           *
*************************************************************************/
Polynomial Polynomial::IntegrateSymbolically(int varIndex) const {
    Polynomial result;
    for (const auto& term : terms) {
        Monomial m = term;
        if (varIndex == 0) { m.px += 1; m.coef /= (double)m.px; }
        else if (varIndex == 1) { m.py += 1; m.coef /= (double)m.py; }
        else if (varIndex == 2) { m.pz += 1; m.coef /= (double)m.pz; }
        result.terms.push_back(m);
    }
    return result;
}
\end{listing}
\normalsize

\pagebreak

\section{Метод Polynomial::Add}
\footnotesize
\begin{listing}{1}
/*************************************************************************
* Метод         : Polynomial::Add                                        *
*------------------------------------------------------------------------*
* Назначение    : Алгебраическое сложение полиномов                      *
* Входные       : other - ссылка на складываемый полином                 *
* Выходные      : Нет (модифицирует текущий объект)                      *
*------------------------------------------------------------------------*
* ОПИСАНИЕ ФУНКЦИОНИРОВАНИЯ:                                             *
* Реализует сложение многочленов с приведением подобных членов.           *
* Алгоритм: 1. Поиск подобных слагаемых по набору степеней (px, py, pz).  *
* 2. Очистка от нулевых членов с коэффициентом < 1e-6 (микропогрешности). *
*************************************************************************/
void Polynomial::Add(const Polynomial& other) {
    for (const auto& otherTerm : other.terms) {
        bool found = false;
        for (auto& myTerm : terms) {
            if (myTerm.px == otherTerm.px && myTerm.py == otherTerm.py && 
                myTerm.pz == otherTerm.pz) {
                myTerm.coef += otherTerm.coef; found = true; break;
            }
        }
        if (!found) terms.push_back(otherTerm);
    }
    terms.erase(remove_if(terms.begin(), terms.end(), 
        [](const Monomial& m) { return abs(m.coef) < 1e-6; }), 
        terms.end());
}
\end{listing}
\normalsize

\pagebreak

\section{Метод Polynomial::Evaluate}
\footnotesize
\begin{listing}{1}
/*************************************************************************
* Метод         : Polynomial::Evaluate                                   *
*------------------------------------------------------------------------*
* Назначение    : Нахождение численного значения полинома в точке        *
* Входные       : x, y, z - координаты расчетной точки                   *
* Выходные      : Вещественное число double - значение в точке           *
*------------------------------------------------------------------------*
* ОПИСАНИЕ ФУНКЦИОНИРОВАНИЯ:                                             *
* Производит численное вычисление значения многочлена в заданной точке   *
* пространства. Математическая модель: вычисление суммы произведений     *
* коэффициентов на переменные в степенях. Используется функция pow       *
* из заголовочного файла <cmath>.                                        *
*************************************************************************/
double Polynomial::Evaluate(double x, double y, double z) const {
    double sum = 0.0;
    for (const auto& t : terms) 
        sum += t.coef * pow(x, t.px) * pow(y, t.py) * pow(z, t.pz);
    return sum;
}
\end{listing}
\normalsize

\pagebreak